(a) If every element of $S_+$ is nilpotent, every prime contains $S_+$ and $\operatorname{Proj}S=\varnothing$. Conversely, if a homogeneous $f\in S_+$ is not nilpotent, then $S_{(f)}$ is a nonzero ring, so $D_+(f)\cong\operatorname{Spec}S_{(f)}$ is nonempty. Thus every positive-degree homogeneous element is nilpotent when $\operatorname{Proj}S=\varnothing$, and so is every finite sum of such elements.

(b) We have $U=\bigcup_{f\in S_+\text{ homogeneous}}D_+(\varphi(f))$, so $U$ is open. Contraction sends its points to homogeneous primes of $S$ avoiding $S_+$. On each indicated open, the graded homomorphism induces $S_{(f)}\to T_{(\varphi(f))}$. The resulting affine morphisms commute with localization and therefore glue to the required morphism.

(c) The opens $D_+(f)$ with $f$ homogeneous of positive degree at least $d_0$ cover $\operatorname{Proj}S$. Their inverse images cover $\operatorname{Proj}T$, since a sufficiently high power of every positive-degree homogeneous element of $T$ is in the image of $\varphi$. In particular, $U=\operatorname{Proj}T$. For each such $f$, the map $S_{(f)}\to T_{(\varphi(f))}$ is an isomorphism: after multiplying numerator and denominator of a fraction by a power of $f$ or $\varphi(f)$, both have degree at least $d_0$, where $\varphi$ is an isomorphism. The local isomorphisms glue.

(d) If $V_i=V\cap D_+(x_i)$, then $A(V_i)=S_{(x_i)}$. Consequently $t(V_i)\cong\operatorname{Spec}S_{(x_i)}\cong D_+(x_i)\subseteq\operatorname{Proj}S$. These identifications agree on overlaps because the restriction maps are the same localizations. They glue to $t(V)\cong\operatorname{Proj}S$.
